%! AP Statistics — Unit 8, Lesson 8.1 — Group Activity
\documentclass[10pt]{article}
\usepackage{apstats-article}
\usepackage{apstats-boxes}

\begin{document}

\pageheader{Unit 8, Lesson 8.1}{Group Activity}
\namepartnerperiod

\vspace{0.1in}
\textbf{Are These Results Surprising?}

\vspace{0.05in}
A school surveys 120 students about their preferred study environment. The results are:

\vspace{0.08in}
\renewcommand{\arraystretch}{1.4}
\begin{tabularx}{0.7\linewidth}{Xcc}
    \textbf{Environment} & \textbf{Observed Count} & \textbf{Expected Count} \\
    \hline
    Library    & 38 & \blank{2cm} \\
    Home       & 52 & \blank{2cm} \\
    Caf\'e     & 18 & \blank{2cm} \\
    Classroom  & 12 & \blank{2cm} \\
    \hline
    \textbf{Total} & \textbf{120} & \blank{2cm} \\
\end{tabularx}

\vspace{0.1in}
\textbf{Directions:} Work with your partner to complete parts (a)--(d) below.

\begin{tcolorbox}[colback=white, colframe=black!40, title=\textbf{Tier R --- Remediate},
    fonttitle=\bfseries, arc=2mm, left=3mm, right=3mm, top=2mm, bottom=2mm]

\textbf{Part (a)} --- Fill in the expected count for each category, assuming all four
study environments are equally preferred.

\vspace{0.05in}
\textit{Hint:} If 120 students are equally spread among 4 categories, expected count
$= 120 \div 4 = $ \blank{2cm} for each category.

\vspace{0.1in}
Go back and fill in the Expected Count column in the table above.
\end{tcolorbox}

\begin{tcolorbox}[colback=white, colframe=black!40, title=\textbf{Tier A --- Approaching Proficiency},
    fonttitle=\bfseries, arc=2mm, left=3mm, right=3mm, top=2mm, bottom=2mm]

\textbf{Part (b)} --- For each category, compute Observed $-$ Expected. Record your
answers below.

\vspace{0.05in}
\begin{tabularx}{\linewidth}{Xcccc}
    & Library & Home & Caf\'e & Classroom \\
    \hline
    Observed $-$ Expected & \blank{1.5cm} & \blank{1.5cm} & \blank{1.5cm} & \blank{1.5cm} \\
\end{tabularx}

\vspace{0.1in}
Which category has the largest \textbf{positive} difference? \blank{3cm}

Which category has the largest \textbf{negative} difference? \blank{3cm}

\vspace{0.1in}
\textbf{Part (c)} --- Without doing any formal calculation, do you think the differences
are large enough to conclude that students do \emph{not} have equal preferences? Explain
your reasoning in 2--3 sentences.

\writelines{3}
\end{tcolorbox}

\begin{tcolorbox}[colback=white, colframe=black!40, title=\textbf{Tier E --- Extension},
    fonttitle=\bfseries, arc=2mm, left=3mm, right=3mm, top=2mm, bottom=2mm]

\textbf{Part (d)} --- What additional information would help you decide whether these
differences are \emph{statistically significant}? In other words, what would a formal
statistical test provide that your informal comparison in part (c) cannot?

\writelines{3}

\textbf{Bonus:} The chi-square statistic adds up all the gaps in a specific way.
Why might we need to \emph{adjust} each gap relative to the expected count, rather than
just adding up the raw gaps?

\writelines{3}
\end{tcolorbox}

\end{document}
